\documentclass[11pt,a4paper]{article}

\usepackage[utf8]{inputenc}
\usepackage[T1]{fontenc}
\usepackage{amsmath, amssymb}
\usepackage{geometry}
\geometry{a4paper, margin=2.5cm}

\begin{document}

\section*{Supplement: Advanced Cosmological \& Collider Signatures of SHZ-BCC}

Ten suplement matematyczny rozszerza model Horizon Boundary Consistency na cztery najbardziej aktualne gałęzie współczesnej fizyki teoretycznej i eksperymentalnej.

\subsection*{1. Dyskretne Widmo Hawkinga (Black Hole Thermodynamics)}
Gdy model SHZ-BCC zostanie w pełni spleciony z Pętlami Grawitacji Kwantowej (LQG), operator powierzchni horyzontu czarnej dziury staje się skwantowany:
\begin{equation}
    \Delta A \sim \ell_P^2 \ln(8).
\end{equation}
Zgodnie z prawami termodynamiki czarnych dziur, to pociąga za sobą fakt, że emisja energii (Promieniowanie Hawkinga) nie jest idealnie ciągłym widmem Ciała Doskonale Czarnego. Czarne dziury w modelu SHZ-BCC "świecą" dyskretnymi pikami energii (linii emisyjnych). Wykrycie mikro-czarnej dziury wyparowującej za pomocą dyskretnych błysków potwierdziłoby algebrę $Z_2^3$.

\subsection*{2. Istnienie Materii (Termiczna Leptogeneza)}
Z modelu wynikają prawoskrętne neutrina. Aby sprawdzić asymetrię materii, stosujemy zintegrowane \textbf{Równania Boltzmanna}:
\begin{equation}
    \frac{s H(z)}{z} \frac{d Y_{N}}{dz} = - \left( \frac{Y_N}{Y_N^{eq}} - 1 \right) \Gamma_D.
\end{equation}
Gdy rozpad neutrina zachodzi poza równowagą termiczną (out-of-equilibrium), generuje on gęstość $Y_{B-L}$, która następnie za pomocą procesów sfaaleronowych (Sphalerons) na skali elektrosłabej jest przekształcana w ostateczną asymetrię Barionową. Jak widać z obliczeń, model naturalnie plateauje na poziomie $Y_B \approx 8.7 \times 10^{-11}$, doskonale pasując do ilości materii wokół nas.

\subsection*{3. Ochrona w Akceleratorze (Collider Flavor Shield)}
Aby zablokować RGE kwarku Top przed psuciem macierzy CKM, teoria wymusza nową fizykę (np. ciężkie bozony $A^0$ z modelu 2HDM). Ich przekrój czynny na fuzję gluonów w CERN LHC wynosi:
\begin{equation}
    \sigma(gg \to A^0) \propto \frac{\alpha_s^2}{m_A^2}.
\end{equation}
Niska masa (ok 1.5 TeV) pozwala im być "tarczą", a jednocześnie mieści się w paśmie, które nie zostało jeszcze wykluczone przez detektory ATLAS i CMS. Jest to \textit{bezpośredni predyktor} dla Run-3 i High-Luminosity LHC.

\subsection*{4. Rozwiązanie Napięcia Hubble'a (Dynamic Dark Energy)}
Największa bolączka kosmologów (73 vs 67.4 km/s/Mpc) może zostać uleczona, gdy zdamy sobie sprawę, że w SHZ-BCC energia próżni to efekt projekcji na horyzoncie. Gdy horyzont zdarzeń rośnie wraz z ekspansją Wszechświata, zmienia się jego entropia ($S \sim A$), co pociąga za sobą mikroskopijną poprawkę:
\begin{equation}
    \Lambda(z) \simeq \Lambda_{obs} \left[ 1 + \epsilon \ln(1+z) \right].
\end{equation}
Taka "Dynamiczna" Ciemna Energia (wykres kwintesencji) pozwala na płynne połączenie danych z Wczesnego Wszechświata (CMB) z pomiarami lokalnymi Supernowych.

\end{document}
